\section{Introduction}
% Stakes
A mobile agent that cooks from a recipe app and then fills a grocery list reads ``vanilla extract'' on screen 1 and must type it on screen 9, after the screen that held the text has left the display. We call a step whose required text appeared on an earlier screen, and only there, a \emph{delayed information need}. Such steps decide whether a long cross-app task completes, and the benchmarks used to develop mobile agents contain few of them. AITZ~\citep{aitz} has none in 18,643 steps, consistent with tasks that average 7.4 steps; GUI-Odyssey~\citep{guiodyssey}, at 15.4 steps per task, has at least one in 31\% of sampled trajectories (\S\ref{sec:needs}).

% Gap
Memory designs for GUI agents miss this need from two sides. Cross-task memories such as AutoDroid~\citep{autodroid} and Agent Workflow Memory~\citep{awm}, and the tips and shortcuts of Mobile-Agent-E~\citep{mobileagente}, store how-to knowledge across episodes and hold no record of what this episode showed ten screens ago. Write-ahead notes and rolling summaries, as in Mobile-Agent-E and GUI-Rise~\citep{guirise}, must guess what will matter before the agent knows, and the needed text sits mid-screen among some 30 elements (\S\ref{sec:needs}). The design that long-document reading suggests is gist-and-lookup~\citep{readagent}: write one short line per screen, let the reader pick a screen from the lines, then re-read that screen in full. Nothing is lost, and the writer can learn to write lines the reader can act on, as document expansion does for search~\citep{doc2query}. We built this \emph{screen index} on GUI-Odyssey and trained the writer from the reader's success at naming a screen that holds the text.

% Finding
The screen index loses to the context it replaces. Full history, the OCR text of every earlier screen, localizes a source screen in 58.7\% of 138 verified cases against 44.2\% for the index (\pval{0.008}) and types the correct text in 34.1\% against 20.3\% (\pval{0.002}). One training round and four prompt variants leave the index within 6 points of its starting point, because one index line has a ceiling. Rewriting the source screen's line alone lifts localization to full-history level (59--66\%) and no higher; an oracle over five such lines gains 27 points only by redundancy that no single writer realizes. Two evaluation choices moved these conclusions by more than any writer effect we measured. Scoring only the \emph{earliest} source screen undercounts the index by 15 points, since 47\% of cases have two or more screens that carry the answer. Keeping cases whose answer is also on the \emph{current} screen (current-screen leakage) turned a 17-point win for our own gating rule into a 23-point loss.

% Contributions
We contribute, on verified and leakage-filtered cases: (1) \textbf{Problem.} A definition, an automatic miner, and the prevalence of delayed information needs (\S\ref{sec:needs}). (2) \textbf{Negative result.} Full history beats a ReadAgent-style screen index on the pooled set (\S\ref{sec:main}). On needs at step 16 or later the lead shrinks to noise; it keeps its localization direction on a 7B reader and on a held-out set, where end-to-end has no winner (\S\ref{sec:length}). (3) \textbf{Why training does not help.} The ceiling of one index line equals full history; the oracle gain is redundancy (\S\ref{sec:writer}). (4) \textbf{Two pitfalls.} Earliest-source scoring and current-screen leakage, with measured effect sizes and the two filter conditions that remove them (\S\ref{sec:pitfalls}).
